% =====================================================================
\section{Running Example and Background}\label{sec:background}
We first introduce the running example, so that every concept that follows can be illustrated on it. We then summarize the concepts the approach builds on: artifact-mediated coordination from MAS research, views and model transformations, incremental rebuilding, and optimistic concurrency control. Table~\ref{tab:glossary} at the end of the section relates each of our terms to its standard counterpart.

\subsection{Running Example: A Document-Service Workflow}\label{sec:motivation}
A small document service is developed by four agents, each responsible for one \emph{view} of the system (Fig.~\ref{fig:network}). The \emph{requirements} view holds service requests and their clauses, such as ``R02: overwriting an archived document returns HTTP status 409 (Conflict)''. The \emph{design} view holds one endpoint per accepted request, with a signature such as \texttt{put(id, body)}. The \emph{implementation} view holds one code patch per endpoint, and the \emph{verification} view holds one check per clause. A shared \emph{security context} holds findings, such as a path-traversal report, that apply to some patches and checks.

The views are heterogeneous: a requirement is neither an operation signature nor a test oracle. Nevertheless, an edit in one view can invalidate artifacts in others. Suppose clause R02 is tightened from ``an archived document cannot be overwritten'' to ``overwriting an archived document returns 409''. The endpoint signature \texttt{put(id, body)} does not change, but the patch must still be regenerated, because the behavior it must implement changed. A dependency analysis that looks only at interfaces would miss this.

\begin{figure}[!t]
\centering
\colorlet{inkblue}{black}\colorlet{pale}{gray!12}\colorlet{resultamber}{gray}% grayscale figure
\begin{tikzpicture}[font=\footnotesize]
\node[box,text width=19mm] (r) at (0,0) {\textbf{Requirements}\\requests, clauses};
\node[box,text width=17mm] (d) at (3.0,0) {\textbf{Design}\\endpoints};
\node[box,text width=22mm] (i) at (6.2,0) {\textbf{Implementation}\\patches};
\node[box,text width=19mm] (t) at (3.0,-1.75) {\textbf{Verification}\\checks};
\node[box,text width=19mm,fill=resultamber!10,draw=resultamber!80!black] (c) at (6.1,-1.75) {\textbf{Shared context}\\security findings};
\draw[arr] (r)--node[above]{$T_1$}(d);
\draw[arr] (d)--node[above]{$T_2$}(i);
\draw[arr] (r.south)|-node[pos=0.25,left]{clauses}(t.west);
\draw[arr] (r.north) to[bend left=22] node[above]{clauses}(i.north);
\draw[arr,dashed,draw=resultamber!80!black] (c)--(i);
\draw[arr,dashed,draw=resultamber!80!black] (c)--(t);
\end{tikzpicture}
\caption{Running example. Solid arrows are declared read dependencies; rule $T_1$ derives an endpoint per request and $T_2$ a patch per endpoint. Dashed arrows carry selected shared context. The controller creates objects and references; agents propose only primitive attribute values (signature text, patch body, check oracle).}\label{fig:network}
\end{figure}

The controller, not an agent, decides \emph{which} artifacts exist: one endpoint per accepted request and one patch per endpoint. Agents fill the \emph{slots} inside those artifacts: an endpoint's signature text, a patch body, a check's oracle. We call one such slot a \emph{job}. Fig.~
ef{lst:job} shows the patch job for \texttt{put} in the form the controller stores it; every element is defined in Section~\ref{sec:semantics}, and the walk-through in Section~\ref{sec:walkthrough} follows this job through a change of R02.

\begin{figure}[!t]
\centering\footnotesize
\renewcommand{\arraystretch}{1.15}
\begin{tabular}{@{}>{\raggedright\arraybackslash}p{0.215\columnwidth}>{\raggedright\arraybackslash}p{0.725\columnwidth}@{}}
\toprule
Element & Value for the patch job of \texttt{put}\\\midrule
Match key & $q=(\textsf{PatchForEndpoint},\ \texttt{endpoint:put})$\\
Job (slot) & $j=(q,\ \texttt{body})$: the body of the patch for \texttt{put}\\
Target & $\iota(q)=\texttt{patch:put}$, created by the controller\\
Owner & \texttt{coding-agent}\\
Read set $F_j$ & signature of \texttt{endpoint:put}; text of clause R02\\
Observation $o_j$ & \textit{source}: signature \texttt{put(id, body)}; R02 = ``overwriting an archived document returns 409''\newline
\textit{context}: SEC-1 = ``path traversal: reject `/'\,''\newline
\textit{contract}: \texttt{impl-v1}; \textit{owner}: \texttt{coding-agent}\\
Key $\Sig_j$ & $\mathrm{enc}(o_j)$: canonical JSON of $o_j$, compared exactly\\
Request $u$ & $(j,\ e{=}7,\ n{=}1)$; records $\Sig_u=\Sig_j$ and $o_u=o_j$ at dispatch\\
Record $B_j$ & after commit: value $v$, validated key $a=\Sig_u$, committing request $\gamma=u$; \emph{usable} while $a=\Sig_j(S)$ and every prerequisite is usable\\
\bottomrule
\end{tabular}
\caption{One job of the running example as the controller stores it. Editing R02 changes $o_j$ and hence $\Sig_j$, so the stored value becomes unusable although the endpoint signature is unchanged, and a reply to $u$ that arrives after the edit cannot commit. Every element is defined in Section~\ref{sec:semantics}.}\label{lst:job}
\end{figure}

\subsection{Coordination Through Shared Artifacts}\label{sec:bg-mas}
MAS research distinguishes \emph{subjective} coordination, in which agents coordinate by reasoning about each other and exchanging messages, from \emph{objective} coordination, in which a coordination medium outside the agents enforces the rules of interaction~\cite{omicini2001coordination}. Linda introduced a shared tuple space as a coordination medium~\cite{gelernter1985linda,gelernter1992coordination}; tuple centres make the medium's reactions programmable, so that coordination laws are specified once, inside the medium, and apply to every agent~\cite{omicini2001tuple}. The A\&A meta-model generalizes the medium to \emph{artifacts}: passive, typed entities that agents use through operations and that guard their own state~\cite{omicini2008artifacts,weyns2007environment}, as realized in CArtAgO and JaCaMo~\cite{ricci2011environment,boissier2013jacamo}. Interaction protocols specify the admissible message exchanges among autonomous agents. Commitment-based protocols give messages a social meaning that can be checked without inspecting the agents~\cite{singh1998acl,yolum2002commitments}, and information-based protocols such as BSPL derive the admissible orderings from which information each message carries, identified by keys~\cite{singh2011bspl,chopra2020protocols}. Electronic institutions add \emph{governors}: runtime mediators that admit an agent's action only if the institution's current state permits it~\cite{esteva2001institutions,arcos2005institutions,boella2006normative}.

We adopt objective, artifact-mediated coordination. The controller is a coordination artifact: agents can only propose values through one operation (reply to a request), and the artifact's guards decide what becomes visible. Section~\ref{sec:related} details what we inherit and what generative proposals add; in short, the added problems are content currentness of generated values, exact and minimal invalidation, and per-episode budgets.

\subsection{LLM Agents as Untrusted Proposers}\label{sec:bg-agents}
In an LLM framework, each \emph{agent} is a model wrapped with a role instruction and possibly with tools~\cite{metagpt,chatdev,autogen}. The framework sends a prompt assembled from current artifacts and receives text, code or JSON. Three properties of this exchange matter. The reply may be wrong in ways that no prompt prevents~\cite{ji2023hallucination}, and the same prompt may yield different replies. A call takes seconds to minutes, so artifacts may change while it is outstanding. Delivery is unreliable: frameworks time out and retry, so a reply may arrive late, twice or never.

\begin{definition}[Request, proposal, validator, contract]\label{def:proposal}
A \emph{request} is one generation call for one job. It has a unique identity and carries the data the agent may see. A \emph{proposal} is a value returned for a request. A \emph{validator} $V_j$ is a deterministic, total predicate that decides whether a proposal is acceptable for the data it was requested on. The \emph{generation contract} $\eta_j$ of job $j$ identifies everything except data that determines what counts as a correct answer: the instruction, output schema, validator and the environment it runs in (for example, tool versions), generation settings, and the owner allowed to answer. Each request records the contract of its job at dispatch. An agent is an \emph{untrusted proposer}: for safety, we assume nothing about its proposals.
\end{definition}

Changing the contract changes the question, just as changing a requirement does: a reply produced under an old instruction or validator must not be reused as if it answered the new one.

\subsection{Views, Transformations and Traces}\label{sec:bg-views}
Model-driven engineering represents each kind of artifact as a \emph{model} that conforms to a \emph{metamodel}, which makes relations between artifacts explicit.

\begin{definition}[Metamodel, view]\label{def:view}
A \emph{metamodel} declares object types with \emph{primitive attributes} (strings, Booleans, numbers or inert code text), typed \emph{references} with multiplicities, and well-formedness predicates. A \emph{view} is a finite graph conforming to a metamodel. Its nodes are objects with stable identities; its edges are typed references.
\end{definition}

\begin{definition}[Rule, match, trace]\label{def:rule}
A transformation \emph{rule} consists of a source pattern, a guard, one target object type, and \emph{bindings} that compute the target's attributes and references. A \emph{match} is a tuple of source objects that fits the pattern and satisfies the guard; each match yields exactly one target object. A \emph{trace} records which match produced which target~\cite{atl}.
\end{definition}

In the running example, rule $\mathsf{RequestToEndpoint}$ creates one endpoint per accepted request with a reference to it; if the request is deleted, the trace identifies the endpoint to delete. Classical transformations compute every binding with an expression. In our setting some bindings are \emph{generated} by an agent: a job is one generated primitive binding of one target. References are always computed by the controller. Consistency checkers re-evaluate rules after model changes~\cite{egyed2006instant,egyed2011tracking}, and views can be pairwise consistent but jointly inconsistent~\cite{stevens,klare}, so a coordinator also needs a global check before release.

\subsection{Incremental Rebuilding and Optimistic Concurrency}\label{sec:bg-incremental}
After a change, a coordinator must decide which jobs to redo; every redone job costs an LLM call. Build systems face the same choice~\cite{mokhov}. Content-based build systems store a key computed from the contents of a task's dependencies and rebuild the task only if the key changes; with \emph{early cutoff}, a rebuilt task that produces the same output as before does not trigger its dependents~\cite{mokhov,mitchell2012shake}. Dependencies may be dynamic~\cite{erdweg2015pluto}. We use the same idea: the \emph{read set} of a job is the source data its agent may read; together with the selected context and the contract it forms the job's \emph{observation}, whose exact encoding is the job's \emph{dependency key}. One difference remains: an LLM is not a deterministic function, so a regenerated value may differ from the old one although both are valid. A value can therefore be checked only by validating it against the observation it was produced from, never by comparing it with a fresh run.

Locking a job's inputs until its agent answers would block every edit for the duration of an LLM call. Databases solve this with \emph{optimistic concurrency control} (OCC)~\cite{occ}: a transaction computes without locks and, at commit time, validates that the data it read has not changed; validation and write are one atomic step. Our commit guard is such a validation phase: a request records the dependency key of the observation it was dispatched on, and a reply commits only if that key still equals the job's current key.

\subsection{Transition Systems and TLA$^+$}\label{sec:bg-tla}
We specify the controller as a transition system. A \emph{safety property} states that nothing bad happens and is proved with an \emph{invariant}; a \emph{liveness property} states that something good eventually happens. TLA$^+$~\cite{lamport1994tla,lamport2002specifying} specifies such systems; the model checker TLC~\cite{yu1999tlc} explores all states of finite instances, and the proof system TLAPS~\cite{cousineau2012tlaps} checks proofs for all instances. A detailed specification \emph{refines} an abstract one if every behavior of the first, mapped through a refinement mapping, is a behavior of the second~\cite{abadi1991refinement,lynch1995simulations}.

\subsection{Coordination Failures and Requirements}\label{sec:failures}\label{sec:requirements}
The running example admits four failure classes that per-answer checks cannot exclude:
\begin{itemize}
\item[\textbf{F1}] \emph{Stale reply.} A request is sent while R02 has revision~1, R02 changes, and the old reply then arrives. It is well formed and may pass the old tests.
\item[\textbf{F2}] \emph{Local validity, global inconsistency.} Design chooses string identifiers and verification integer identifiers. Both pass their local validators.
\item[\textbf{F3}] \emph{Unsafe change of team or rules.} A reviewer is added, a permission is revoked, or a prompt or validator changes, and old evidence or old completion verdicts are reused as if nothing had changed.
\item[\textbf{F4}] \emph{Unbounded or wasted generation.} A failing request is retried on every scheduler pass, or an unrelated edit regenerates every artifact.
\end{itemize}
A coordinator for untrusted generative agents should therefore satisfy six requirements; Section~\ref{sec:guarantees} answers each with a theorem.
\begin{itemize}
\item[\textbf{R1}] \emph{Structural integrity.} Objects, traces and references are well typed and have exactly one producer, whatever the agents return.
\item[\textbf{R2}] \emph{Currentness.} A value is usable only if it was validated against the \emph{current} observation, contract and owner; obsolete and duplicate replies cannot write (F1, F3).
\item[\textbf{R3}] \emph{Honest completion.} Publication requires that nothing is pending or stale and that all joint constraints hold (F2, F3).
\item[\textbf{R4}] \emph{Locality.} A change invalidates exactly the work whose observations changed (F4).
\item[\textbf{R5}] \emph{Bounded effort and termination.} Generation attempts are bounded, and a stable workflow ends with an explicit success or failure verdict (F4).
\item[\textbf{R6}] \emph{Safe reconfiguration.} Changing agents, rules, permissions or validators preserves R1--R3 (F3).
\end{itemize}

\begin{table*}[!t]
\centering
\caption{Terms used in this paper, their symbols, and their standard counterparts in MAS coordination, databases and build systems. Formal definitions are in Section~\ref{sec:semantics}.}\label{tab:glossary}
\footnotesize\setlength{\tabcolsep}{3pt}
\begin{tabular}{@{}P{23mm}P{14mm}P{68mm}P{62mm}@{}}\toprule
Term & Symbol & Meaning in this paper & Standard counterpart\\\midrule
Controller & -- & Deterministic component that owns structure, keys, acceptance and publication & Coordination artifact (A\&A); tuple centre; institutional governor\\
Job (slot) & $j=(q,b)$ & One generated primitive attribute $b$ of the target of match $q$ & Build task; artifact operation with one output\\
Read set & $F_j$ & Source data the job's agent may read (a pure query, including membership) & Read set of a transaction; declared/dynamic build dependencies\\
Contract & $\eta_j$ & Instruction, schema, validator and its environment, generation settings, owner & Build rule plus toolchain version; protocol role\\
Observation & $o_j$ & Read set, selected context and contract: all the job depends on & Materialized read set plus task configuration\\
Dependency key & $\Sig_j$ & Exact injective encoding $\operatorname{enc}(o_j)$ & Content-addressed rebuild key (verifying trace)\\
Validated key & $a_j$ & The key against which the stored value was validated & Read-set version validated at OCC commit\\
Episode & $e_j$, $n_j$ & Maximal period in which a job keeps one key; attempt counter $n_j\le k$ & Invalidation epoch with a retry budget\\
Request & $u=(j,e,n)$ & One generation call; identity never reused & Message identifier / BSPL key; transaction identifier\\
Accepted / usable & $U_j$ & Accepted: a value was once validly committed. Usable: its validated key is current and all prerequisites are usable & Committed vs.\ current version; up-to-date build target\\
Publishable & $\mathit{Pub}$ & Nothing suspended, stale or pending; global predicate $G$ holds (Eq.~\eqref{eq:publish}, written $\Pub(S)$) & Release gate; institutional completion condition\\
Suspended stage & $Q$ & Stage whose prerequisites are unavailable; contributes no jobs & Blocked task; disabled scene in an institution\\\bottomrule
\end{tabular}
\end{table*}
